\chapter{Einleitung}\label{chap:Einleitung}
Moderne \gls{ic} bestehen aus einer Vielzahl digitaler und analoger Komponenten, die nach dem Einschalten der Versorgungsspannung einen definierten Anfangszustand benötigen, um fehlerfrei zu funktionieren. Besonders bei komplexeren Systemen, in denen Speicherzellen, Flip-Flops oder digitale Steuerlogik zum Einsatz kommen, ist eine saubere Initialisierung essenziell. Ohne eine gezielte Zurücksetzung können sich undefinierte Betriebszustände einstellen, die die Funktionalität des gesamten Systems gefährden.

Um eine definierte Initialisierung sicherzustellen, wird eine \gls{por}-Schaltung eingesetzt. Diese überwacht die anliegende Versorgungsspannung und erzeugt beim Einschalten ein internes Reset-Signal, welches die relevanten Komponenten des \glspl{ic} in einen vordefinierten Ausgangszustand versetzt. Gleichzeitig stellt die \gls{por}-Schaltung sicher, dass bei Spannungsabfällen oder instabilen Betriebsbedingungen ein Reset ausgelöst wird, um Fehlfunktionen zu verhindern.

Im Rahmen dieser Arbeit wird eine \gls{por}-Schaltung für einen neu entwickelten Chip entworfen. Der Chip beinhaltet sowohl analoge als auch digitale Funktionsblöcke, weshalb separate \gls{por}-Schaltungen für den analogen als auch digitalen Bereich entwickelt werden. Aufgrund unterschiedlicher Versorgungsspannungen und technologischer Anforderungen der beiden Bereiche müssen die Schaltungen an die jeweiligen Bedingungen angepasst werden.

Die Anforderungen an zuverlässige \gls{por}-Schaltungen steigen insbesondere bei modernen Fertigungstechnologien wie dem hier betrachteten 130-nm-CMOS-Prozess, da sich durch Prozessschwankungen, Temperatureinflüsse und Versorgungsspannungsänderungen das Schaltverhalten der Transistoren und passiven Bauelemente erheblich verändern kann. Ziel dieser Arbeit ist die Entwicklung und Verifizierung einer robusten \glspl{por}-Schaltung, die auch unter variierenden \gls{pvt} zuverlässig arbeitet. Dabei liegt der Fokus auf einer stabilen Referenzspannung und einer zuverlässigen Reset-Verzögerung.

Durch Simulationen in der Cadence-Virtuoso-Umgebung wird sichergestellt, dass die entwickelte \gls{por}-Schaltung den vorgegebenen Anforderungen entspricht und eine zuverlässige Initialisierung des analogen und des digitalen Chipbereichs gewährleistet wird.
